\documentclass[11pt,letterpaper]{article}
\usepackage[T1]{fontenc}
\usepackage{libertinus}
\usepackage[margin=1in,headheight=15pt]{geometry}
\usepackage{amsmath,amssymb,amsthm,mathtools}
\usepackage{microtype,booktabs,tabularx,longtable,enumitem}
\usepackage{xcolor,xurl}
\usepackage[hidelinks]{hyperref}
\usepackage[nameinlink,noabbrev]{cleveref}
\usepackage{fancyhdr}
\pagestyle{fancy}\fancyhf{}
\fancyhead[L]{\small Gaussian dust and Christoffel recovery}
\fancyhead[R]{\small ProveIt research extension}
\fancyfoot[C]{\thepage}
\setlength{\parskip}{0.35em}
\setlength{\parindent}{1.2em}
\setlist{itemsep=0.3em,topsep=0.4em}
\emergencystretch=2em
\numberwithin{equation}{section}
\newtheorem{theorem}{Theorem}[section]
\newtheorem{proposition}[theorem]{Proposition}
\newtheorem{lemma}[theorem]{Lemma}
\newtheorem{corollary}[theorem]{Corollary}
\theoremstyle{definition}
\newtheorem{definition}[theorem]{Definition}
\newtheorem{question}{Research question}
\theoremstyle{remark}
\newtheorem{remark}[theorem]{Remark}
\newcommand{\R}{\mathbb R}
\newcommand{\N}{\mathbb N}
\newcommand{\E}{\mathbb E}
\newcommand{\Prob}{\mathbb P}
\newcommand{\TV}{\operatorname{TV}}
\newcommand{\Law}{\mathcal L}
\newcommand{\Var}{\operatorname{Var}}
\newcommand{\sinc}{\operatorname{sinc}}
\newcommand{\supp}{\operatorname{supp}}
\newcommand{\norm}[1]{\left\lVert#1\right\rVert}
\newcommand{\dd}{\,\mathrm d}
\newcommand{\K}{\mathcal K_V}
\newcommand{\Cs}{\mathcal C}
\newcommand{\clip}{\operatorname{clip}_{[0,V]}}
\newcommand{\pin}{458ccfb80b9940220819091da107e82bb970dcca}
\hypersetup{pdftitle={Gaussian Dust and Christoffel Recovery in Infinite Uniform Convolutions},pdfauthor={AI-assisted research manuscript prepared for Vladimir Reshetnikov},pdfsubject={Fabius--Rvachev inverse problems, Gaussian variance, moment recovery}}
\title{\textbf{Gaussian Dust and Christoffel Recovery\\in Infinite Uniform Convolutions}\\[0.55em]\large Exact non-estimability, a complete compact-class criterion,\\and geometric moment formulas beyond finite capacity}
\author{Research manuscript prepared for Vladimir Reshetnikov\\\small AI-assisted; conventional proofs and reproducible exact checks}
\date{29 September 2026}
\begin{document}
\maketitle
\begin{abstract}
The finite-capacity inverse problems in ProveIt's Fabius--Rvachev research
program change qualitatively when the number of uniform convolution factors
is unrestricted. We study a known smooth compactly supported background plus
an independent Gaussian and any square-summable sequence of centered uniform
factors. The underlying distributional compactification and identifiability
belong to the existing DUSTPAN theory of Billey--Swanson; they are credited and
re-established in a variance-weighted moment representation suited to inference.
In this experiment the Gaussian variance is identifiable, but its exact
minimax absolute-error risk is $V/2$, and its exact squared-error risk is
$V^2/4$, at every finite sample size when total latent variance is at most $V$.
On a compact subclass, uniform consistency is possible if and only if the
uniform-factor energy tails vanish uniformly. We provide finite-cumulant
upper certificates, a Christoffel hierarchy with an explicit tail bound,
and a rank-free total-variation expansion quantifying Gaussian approximation.
For squared scales $Cr^j$, the degree-$d$ Christoffel bias is exactly
$C(1-r)r^d/[3(1-r^{d+1})^2]$. Coefficient-constrained regularization gives
noise bounds and statistical upper rates on two-sided polynomial and geometric
envelope classes. The article distinguishes classical tools, derived results,
and unresolved minimax and computational questions. No proof-assistant
verification or worldwide priority claim is made.
\end{abstract}
\vspace{0.4em}
\noindent\textbf{Keywords.} Fabius function; Rvachev up-function; DUSTPAN distributions;
Gaussian component; infinite convolution; Christoffel function; cumulants;
statistical identifiability; uniform consistency.
\clearpage
\tableofcontents
\newpage

\section{Research question, provenance, and main conclusions}
\label{sec:intro}

\subsection{What is being extended}
ProveIt contains exact Fabius--Rvachev arithmetic and several recent inverse
problem reports. The finite-factor stability report studies a fixed number of
unknown uniform factors against a known smooth background. Its successor
allows an unknown Gaussian variance but keeps the capacity fixed. The earlier
infinite-spectrum report constructs moment-preserving perturbations and
quantitative inverse-instability examples near the dyadic spectrum
\cite{repo,finite,gaussian,recovering}. These are different questions: neither
finite-capacity rates nor full-law identifiability alone determine what happens
to Gaussian-variance estimation when arbitrarily many tiny factors are allowed.

We address the following specific extension:
\begin{quote}
For an unrestricted square-summable uniform spectrum, can one distinguish
actual Gaussian variance from the accumulated variance of unresolved small
uniform factors? Which restrictions make this distinction statistically
recoverable, and can the recovery be certified from finitely many cumulants?
\end{quote}
The answer has both a negative and a constructive part. The unrestricted
statistical problem has exactly the risk of ignoring the data. On compact
subclasses, one precise tail condition is necessary and sufficient for uniform
consistency. Positive polynomial certificates turn this qualitative condition
into an explicit method.

\subsection{Classical results that are not claimed as new}
Billey and Swanson introduced DUSTPAN distributions, characterized their
compact parameter space, and proved the closure and identifiability statements
for generalized uniform sums plus a normal component
\cite[Theorem~1.15, Theorem~3.13, Corollaries~3.28--3.29]{billey}.
Their interval lengths are twice our half-lengths. Thus the existence of a
Gaussian component in the closure, the uniqueness of the spectrum, and the
basic compactification are \emph{not} discoveries of this article.

The variational definition and moment-matrix formula for Christoffel functions
are also classical; \cite[Section~2.2]{christoffel} gives a modern account,
including singular moment matrices. Our geometric computation is an elementary
Cauchy-determinant specialization, related to the geometric orthogonality
measures in the little $q$-Jacobi family \cite[Section~18.27(iv)]{dlmfq}.
No novelty is asserted for those identities or for general compact moment
determinacy. The distinction between exact moment determination and unstable
inversion has an extensive literature; \cite{hausdorff} treats a different,
operator-theoretic Hausdorff moment setting.

ProveIt also already contains Christoffel reconstruction of the physical
measure $\mathrm{up}(x)\dd x$, as recorded in its representation-frontier
corpus audit \cite{native}. Our measure is instead
$3s\delta_0+\sum_j a_j^2\delta_{a_j^2}$: its atoms encode convolution
factors, not spatial observations of the up-density. The distinction prevents
conflating an existing density-reconstruction result with the variance-recovery
problem treated here.

The contribution developed here is the combination of a variance-weighted
spectral measure with the smoothed statistical experiment: exact finite-sample
risk, the compact-class equivalence theorem, explicit rank-free $L^1$
remainder bounds, and constructive Gaussian-variance certificates with exact
geometric bias and controlled noisy implementations. These results have full
proofs below. Their worldwide publication priority has not been established.

\subsection{A guide to the theorems}
Fix a variance budget $V>0$. In the notation introduced in \cref{sec:model},
$s$ is Gaussian variance, $v$ is \emph{total latent variance}, and $a_j$ are
uniform half-lengths.
\begin{center}
\begin{tabularx}{\textwidth}{@{}p{0.27\textwidth}X@{}}
\toprule
Question & Answer proved here or credited below \\\midrule
What does the exact law identify? & All $a_j$, with multiplicity, and $s$; the
underlying identification is classical. The spectral measure in
\cref{thm:spectral} makes the Gaussian variance an atom at zero. \\
What can a finite sample learn without tail restrictions? & Exactly minimax
risk $V/2$ in absolute error and $V^2/4$ in squared error,
\cref{thm:minimax}. \\
Which compact classes are learnable? & Exactly those with uniformly vanishing
$\ell^2$ tails of $a$, \cref{thm:criterion}. \\
Can finite cumulants certify recovery? & Yes: decreasing positive certificates
and a Christoffel hierarchy, \cref{thm:ladder,thm:christoffel}. \\
What happens for geometric factors? & An exact finite-degree bias formula,
\cref{thm:geometric}; the dyadic specialization is displayed explicitly. \\
How rapidly can tiny factors imitate Gaussian noise? & A total-variation
asymptotic with an explicit uniform remainder, \cref{thm:dust,cor:dust}. \\
\bottomrule
\end{tabularx}
\end{center}

\subsection{Snapshot and verification boundary}
The inspected repository snapshot is commit
\begin{center}\small\texttt{\pin}.\end{center}
The relevant directory is
\begin{quote}\small\path{Analysis/FabiusFunction/docs/semi-formalized-research-frontiers/drafts/inverse-and-sampling/}.\end{quote}
The accompanying source audit records the particular files inspected; this is
not an audit of every theorem in ProveIt. None of the proofs below assumes
that an unreviewed repository theorem is correct. Supporting computations
check finite algebraic identities, not the universal analytic or statistical
claims. The manuscript has not been independently refereed or checked in
Lean or Rocq. No repository files were changed.

\section{Model and elementary analytic facts}
\label{sec:model}

Let $U_1,U_2,\ldots$ be independent uniform random variables on $[-1,1]$,
and let $G$ be an independent standard normal. Let $B$ be independent of them,
known, centered, and supported in $[-H,H]$, with density
$b\in C_c^\infty(\R)$. Centering a known background causes no loss of generality.
The statistical observations are independent copies of
\begin{equation}
 Y=B+Z,\qquad Z=\sqrt{s}\,G+\sum_{j\ge1}a_jU_j,
 \qquad a_1\ge a_2\ge\cdots\ge0,
 \quad \sum_{j\ge1}a_j^2<\infty.
 \label{eq:model}
\end{equation}
The series converges in $L^2$ and almost surely, by independence, centering,
and the summability of the variances. The $L^2$ statement follows directly
from the Cauchy criterion; the almost-sure statement is Kolmogorov's
convergence theorem, also recalled in \cite[Theorem~3.9]{billey}.
Only the $L^2$ realization is needed in the subsequent arguments.

Write
\begin{equation}
 x_j=a_j^2,\qquad p_k(a)=\sum_{j\ge1}x_j^k,\qquad
 v=s+\frac{p_1(a)}3,\qquad 0\le v\le V,
 \qquad L=3V.
 \label{eq:notation}
\end{equation}
Thus $s$ and $v$ are variances, not standard deviations. The latent law is
$\mu_{v,a}=\Law(Z)$ and the observed law is $P_{v,a}=\Law(Y)$.
We use $\TV(P,Q)=\sup_A|P(A)-Q(A)|$, equal to half the $L^1$ distance of
densities. Gaussian measure of variance $t\ge0$ is denoted by $\gamma_t$;
$\gamma_0=\delta_0$.

\begin{lemma}[Uniform exponential and moment bounds]
\label{lem:subgaussian}
For every real $t$,
\begin{equation}
 \E e^{tZ}\le e^{vt^2/2}\le e^{Vt^2/2}.
 \label{eq:subgaussian}
\end{equation}
The moment-generating and characteristic functions of $Z$ extend to entire
functions. For a constant $D$ depending only on $H,V$, and every integer
$r\ge1$, $\E|Y|^r\le(D\sqrt r)^r$ uniformly over the model.
\end{lemma}
\begin{proof}
The coefficient inequality $(2k+1)!\ge6^k k!$ follows by induction because
$(2k+3)(2k+2)\ge6(k+1)$. Therefore
$\sinh t/t\le e^{t^2/6}$. Apply this to each finite partial sum and to its
independent Gaussian. The partial sums converge in probability, and their
exponentials are uniformly integrable: their second exponential moments are
bounded by the same inequality at $2t$. Passing to the limit proves
\eqref{eq:subgaussian} and the product representation of the MGF.

The bound also holds with $t$ and $-t$, so
$\E e^{R|Z|}\le2e^{VR^2/2}$ for all $R$. Differentiation under expectation
on compact subsets of $\mathbb C$ proves entire extension. Chernoff's
inequality gives $\Prob(|Z|>u)\le2e^{-u^2/(2V)}$. Integration of this tail
bound yields $\E|Z|^r\le(D_0\sqrt r)^r$ after increasing a constant $D_0$.
The triangle inequality in $L^r$ and $|B|\le H$ give the assertion for $Y$.
\end{proof}

\begin{lemma}[Cumulant coordinates]
\label{lem:cumulants}
Set $\gamma_{2k}=\kappa_{2k}(U_1)$. Then
\begin{equation}
 \gamma_{2k}=\frac{2^{2k}B_{2k}}{2k}
 =(-1)^{k+1}\frac{2(2k-1)!\zeta(2k)}{\pi^{2k}}\ne0,
 \label{eq:gamma}
\end{equation}
and
\begin{align}
 \kappa_2(Y)-\kappa_2(B)&=v,\label{eq:variance}\\
 \kappa_{2k}(Y)-\kappa_{2k}(B)&=\gamma_{2k}p_k(a),\quad k\ge2.
 \label{eq:cumulants}
\end{align}
The odd latent cumulants vanish. In particular,
$\gamma_2=1/3$, $\gamma_4=-2/15$, and $\gamma_6=16/63$.
\end{lemma}
\begin{proof}
Expand $\log(\sinh t/t)$ at zero using the classical logarithmic sine
expansion \cite[Eq.~4.19.7]{dlmf}. Equivalently, Euler's sine product gives
$\log(\sinh t/t)=\sum_{k\ge1}(-1)^{k+1}\zeta(2k)t^{2k}/(k\pi^{2k})$
for $|t|<\pi$. This proves both forms of \eqref{eq:gamma}.
For sufficiently small $|t|$, the sum of these logarithmic series over $j$
converges absolutely, since $\sum a_j^2<\infty$ and the $a_j$ are bounded.
Comparing coefficients of the independent-product MGF proves the formulas.
\end{proof}

\begin{remark}[The Rvachev background]
The random variable $\sum_{j\ge1}2^{-j}U_j$ has the normalized Rvachev
up density on $[-1,1]$. Its characteristic function is
$\prod_{j\ge1}\sinc(2^{-j}t)$ in our ordinary characteristic-function
convention, and its variance is $1/9$. It is an admissible known background
in \eqref{eq:model}, as in the inspected Gaussian-confounding report
\cite{gaussian}. Alternatively, the same geometric spectrum can be the
\emph{unknown latent} uniform component. These two roles must not be confused.
\end{remark}

\section{A spectral measure for the Gaussian component}
\label{sec:spectral}

\begin{definition}[Variance-weighted spectral measure]
For the parameters in \eqref{eq:notation}, define a finite measure on $[0,L]$ by
\begin{equation}
 \rho_{v,a}=3s\,\delta_0+\sum_{j\ge1}x_j\delta_{x_j}.
 \label{eq:spectral}
\end{equation}
Zero coordinates contribute no mass. At any positive $x$, the mass is $x$
times the finite integer multiplicity of that squared scale.
\end{definition}

\begin{theorem}[Moment representation and exact identification]
\label{thm:spectral}
The moments $m_r=\int x^r\dd\rho_{v,a}(x)$ are
\begin{equation}
 m_0=3v,\qquad m_r=p_{r+1}(a)\quad(r\ge1).
 \label{eq:spectralmoments}
\end{equation}
Consequently the exact observed law determines $\rho_{v,a}$, all the positive
half-lengths with multiplicity, and $s=\rho_{v,a}(\{0\})/3$.
\end{theorem}
\begin{proof}
The mass identity is $3s+p_1=3v$; all positive-degree moments kill the atom
at zero and give \eqref{eq:spectralmoments}. The cumulant formulas recover
these moments from the observed law after subtracting the known background
cumulants. Two finite measures on $[0,L]$ with the same moments integrate
every polynomial equally. Uniform polynomial approximation of continuous
functions, followed by uniqueness of finite Borel measures, makes them equal.
At each positive atom $x$, divide its mass by $x$ to recover its multiplicity,
and take its positive square root to recover the half-length. The atom at
zero gives $s$.
\end{proof}

This is a moment-theoretic representation of the classical identification
result, not a new claim that generalized uniform convolutions are identifiable.
Its additional usefulness is that Gaussian variance is now a boundary-atom
estimation problem for a positive measure whose entire moment sequence is
known in principle.

\begin{proposition}[Intrinsic Gaussian divisibility]
\label{prop:intrinsic}
For every parameter in the model,
\begin{equation}
 s=\lim_{t\to+\infty}\frac{2\log\E e^{tY}}{t^2}
   =\max\{w\ge0:P_{v,a}=\gamma_w*Q\text{ for a probability measure }Q\}.
 \label{eq:intrinsic}
\end{equation}
\end{proposition}
\begin{proof}
For each fixed $a_j$, $t^{-2}\log(\sinh(a_jt)/(a_jt))\to0$; it is
nonnegative and bounded by $a_j^2/6$. Dominated convergence for the sum
shows that the uniform component contributes $o(t^2)$ to the log MGF.
The bounded background contributes at most $H|t|$ in absolute value.
The Gaussian contributes $st^2/2$, proving the limit.

The decomposition with $w=s$ exists by construction. If $P_{v,a}=\gamma_w*Q$,
Tonelli's theorem applied to exponential moments shows that $Q$ has finite
MGF at every real $t$ and
$\log M_Q(t)=\log M_Y(t)-wt^2/2$.
Its mean is zero. Jensen's inequality gives $\log M_Q(t)\ge0$.
The displayed limit therefore implies $w\le s$.
\end{proof}

\section{The known compactification and its statistical topology}
\label{sec:topology}

Define
\begin{equation}
 \K=\left\{(v,a):0\le v\le V,\quad
 a_1\ge a_2\ge\cdots\ge0,\quad\sum_{j\ge1}a_j^2\le3v\right\},
 \label{eq:K}
\end{equation}
with metric
$d_\star((v,a),(w,c))=|v-w|+\sup_j|a_j-c_j|$.
Here $s(v,a)=v-p_1(a)/3$. The point of retaining $v$, rather than $s$, as a
coordinate is that variance carried by disappearing factors is not lost.
The following proposition supplies a self-contained smoothed version of the
Billey--Swanson framework \cite{billey}.

\begin{proposition}[Compact model and a total-variation homeomorphism]
\label{prop:compact}
The space $\K$ is compact. Its map to latent laws is continuous and injective
for weak convergence, and its map to observed laws is a homeomorphism onto a
compact subset of the total-variation metric space. Finite uniform sums with
$s=0$ are dense in this observed-law space.
\end{proposition}
\begin{proof}
Monotonicity gives $a_j^2\le L/j$. A diagonal subsequence of any parameter
sequence converges in each coordinate and in $v$; Fatou's lemma preserves
the defining inequality. The uniform tail bound $a_j\le\sqrt{L/j}$ upgrades
coordinate convergence to convergence in the supremum norm. This proves
compactness.

For a complex argument $z$, write the compensated product
\begin{equation}
 \phi_{v,a}(z)=e^{-vz^2/2}
 \prod_{j\ge1}\left(e^{a_j^2z^2/6}\sinc(a_jz)\right).
 \label{eq:compensated}
\end{equation}
On every compact set in $z$, a factor is $1+O(a_j^4)$ uniformly for large $j$,
and
\begin{equation}
 \sum_{j>J}a_j^4\le a_{J+1}^2\sum_{j>J}a_j^2\le\frac{L^2}{J+1}.
 \label{eq:tailfour}
\end{equation}
Finite products depend continuously on the coordinates, and the tails converge
uniformly to one as $J\to\infty$. Hence the entire characteristic functions
converge locally uniformly when parameters converge. L\'evy's continuity
theorem gives weak continuity. Injectivity is \cref{thm:spectral}, even
without the background.

If $\mu_n\Rightarrow\mu$, then for every $y$,
$\int b(y-z)\dd\mu_n(z)\to\int b(y-z)\dd\mu(z)$, because $b(y-\cdot)$ is
bounded and continuous. These nonnegative densities all have integral one.
Scheff\'e's lemma gives $L^1$ convergence: indeed, dominated convergence for
$\min(f_n,f)\le f$ and
$\norm{f_n-f}_1=2-2\int\min(f_n,f)$ proves it directly.
Thus the observed-law map is TV-continuous. It is injective by
\cref{thm:spectral}, and a continuous injection from a compact space into a
metric space is a homeomorphism onto its image.

For density, retain the first $J$ half-lengths and put
$\tau_J=v-\frac13\sum_{j\le J}a_j^2\ge0$.
Replace this residual variance by $n_J$ uniforms of half-length
$\sqrt{3\tau_J/n_J}$, with $n_J\to\infty$. Choose $n_J$ so that this
half-length tends to zero, and sort the combined finite list. Its total
variance is $v$, its Gaussian variance is zero, and each ordered coordinate
tends to the target coordinate. The same uniform-tail argument yields
$d_\star$ convergence, hence TV convergence of the observed laws.
\end{proof}

\begin{theorem}[Exactly where Gaussian variance is continuous]
\label{thm:continuity}
The function $s:\K\to[0,V]$ is upper semicontinuous. Its continuity points
are exactly the parameters with $s=0$. At a parameter with Gaussian variance
$s$, the possible limits of $s(\theta_n)$ along sequences
$\theta_n\to\theta$ are precisely the interval $[0,s]$.
\end{theorem}
\begin{proof}
The map $p_1$ is the increasing supremum of continuous finite sums, so it is
lower semicontinuous; $s=v-p_1/3$ is upper semicontinuous. At $s=0$,
nonnegativity and upper semicontinuity imply continuity.

Fix a target and a number $u\in[0,s]$. Keep the original uniform spectrum,
replace Gaussian variance $s-u$ by $n$ equal uniform factors of half-length
$\sqrt{3(s-u)/n}$, and keep Gaussian variance $u$. Sorting changes no law.
The total variance stays $v$, and the ordered coordinates converge to the
original spectrum. The Gaussian variances are constantly $u$. This proves
attainability of every proposed limit and discontinuity when $s>0$.
Upper semicontinuity and nonnegativity exclude all other limits.
\end{proof}

\begin{remark}
The full parameter $(v,a)$ is uniformly recoverable as a \emph{continuous
function of the exact law} in the compactified metric. The functional $s$ is
not continuous on that parameter space. There is no contradiction: infinitely
many small coordinate changes may carry a fixed sum of squares.
At a zero-Gaussian target, convergence also holds in $\ell^2$ for the spectra,
since their squared norms converge and their coordinates converge.
\end{remark}

\section{A rank-free expansion for Gaussian dust}
\label{sec:dust}

The compactification describes convergence but does not give a numerical error.
The next theorem quantifies how a cloud of small uniform factors replaces a
Gaussian component. The background in this section need only be a probability
density $b\in W^{8,1}(\R)$; compact support is not required. Write
$\norm{b^{(r)}}_1$ for the $L^1$ norm of its weak derivative.

\begin{theorem}[A uniform fourth-order smoothing expansion]
\label{thm:dust}
Let $a$ be any finite or square-summable sequence of nonnegative half-lengths,
and put
\[
 f=b*\Law\left(\sum_j a_jU_j\right),\qquad
 g=b*\gamma_{p_1/3}.
\]
Then
\begin{equation}
 \left\|f-g+\frac{p_2}{180}g^{(4)}\right\|_1
 \le \frac{43}{22680}p_3\norm{b^{(6)}}_1
      +\frac{p_2^2}{4050}\norm{b^{(8)}}_1.
 \label{eq:dustbound}
\end{equation}
The constants are independent of the number of factors and of their collisions.
They are not asserted optimal.
\end{theorem}
\begin{proof}
Let $T_i$ be convolution by $a_iU_i$, let $G_i$ be convolution by
$\gamma_{a_i^2/3}$, and set $A_i=T_i-G_i$. All these operators commute with
weak derivatives, and $T_i,G_i$ contract $L^1$. Taylor's formula for
translation, with integral remainder, gives
\begin{equation}
 \left\|h(\cdot-z)-\sum_{r=0}^{k-1}
          \frac{(-z)^r}{r!}h^{(r)}\right\|_1
 \le\frac{|z|^k}{k!}\norm{h^{(k)}}_1.
 \label{eq:translationbound}
\end{equation}
It applies to Sobolev functions by approximation from smooth functions.
The uniform and its variance-matched Gaussian agree in moments through order
three. Their fourth moments are $a_i^4/5$ and $a_i^4/3$, respectively.
Consequently
\begin{align}
 \norm{A_i h}_1&\le\frac{a_i^4}{45}\norm{h^{(4)}}_1,
 \label{eq:Ai4}\\
 \left\|A_i h+\frac{a_i^4}{180}h^{(4)}\right\|_1
 &\le\frac{11a_i^6}{11340}\norm{h^{(6)}}_1.
 \label{eq:Ai6}
\end{align}
For the first line use $(1/5+1/3)/4!=1/45$.
For the second line expand through degree five, use the fourth-moment
\emph{difference} $(1/5-1/3)/4!=-1/180$, and bound the sixth-moment
remainder by $(1/7+5/9)/6!=11/11340$.

First suppose there are $n$ factors. An exact telescoping identity is
\[
 \prod_{i=1}^n T_i-\prod_{i=1}^n G_i
 =\sum_{i=1}^n A_i
       \left(\prod_{j<i}T_j\right)\left(\prod_{j>i}G_j\right).
\]
Replace the product $\prod_{j<i}T_j$ by $\prod_{j<i}G_j$ with a second
telescoping identity. Every resulting remainder contains $A_iA_j$ for
one pair $j<i$, and all other factors are $L^1$ contractions. Applying
\eqref{eq:Ai4} twice bounds its action on $b$ by
$a_i^4a_j^4\norm{b^{(8)}}_1/45^2$. Thus
\begin{equation}
 f-g=\sum_i A_i\left(\prod_{j\ne i}G_j\right)b+R,
 \qquad
 \norm R_1\le\frac{p_2^2}{4050}\norm{b^{(8)}}_1.
 \label{eq:doubletelescoping}
\end{equation}
There are no uncontrolled higher-order subset sums in this argument.

Apply \eqref{eq:Ai6} to each summand. To replace
$\left(\prod_{j\ne i}G_j\right)b^{(4)}$ by $g^{(4)}$, use centering and
\eqref{eq:translationbound} at order two:
\[
 \norm{(I-G_i)h}_1\le\frac{a_i^2}{6}\norm{h''}_1.
\]
Its contribution, after multiplication by $a_i^4/180$, is at most
$a_i^6\norm{b^{(6)}}_1/1080$. Since
\[
 \frac{11}{11340}+\frac1{1080}=\frac{43}{22680},
\]
we obtain \eqref{eq:dustbound} for a finite list.

For an infinite list, the uniform partial sums converge weakly, and convolution
with any $L^1$ function is continuous from weak convergence into $L^1$.
For completeness, approximate that function in $L^1$ by a continuous
compactly supported function; for the latter, tightness controls spatial tails
and dominated convergence applies on a compact interval. Convolution
contractions control the approximation errors. Apply this observation to
$b$ and its derivatives. Gaussian convolutions have the same continuity,
and the finite $p_1,p_2,p_3$ sums converge to their infinite counterparts.
Pass to the limit in the finite-list inequality.
\end{proof}

\begin{corollary}[Sharp first response of a small-factor cloud]
\label{cor:dust}
For a sequence of spectra with $h=\sup_j a_j\to0$,
$p_1\to3\tau<\infty$, and $p_2>0$, one has
\begin{equation}
 \TV\left(b*\Law\left(\sum_j a_jU_j\right),
               b*\gamma_{p_1/3}\right)
 \sim\frac{p_2}{360}\norm{(b*\gamma_\tau)^{(4)}}_1.
 \label{eq:dustasymptotic}
\end{equation}
The constant multiplying $p_2$ is strictly positive.
\end{corollary}
\begin{proof}
The inequalities $p_3\le h^2p_2$ and $p_2\le h^2p_1$ make the remainder in
\eqref{eq:dustbound} equal to $o(p_2)$. Moreover,
$(b*\gamma_{p_1/3})^{(4)}\to(b*\gamma_\tau)^{(4)}$ in $L^1$.
If the latter derivative were zero, the density would be a polynomial of
degree at most three in the distributional sense. An integrable polynomial
on $\R$ must be zero, contradicting mass one. Divide the expansion by $p_2$,
take $L^1$ norms, and divide by two.
\end{proof}

\paragraph{Equal dust.}
Take $n$ factors with $a_j=\sqrt{3\delta/n}$. They have variance $\delta$,
$p_2=9\delta^2/n$, and $p_3=27\delta^3/n^2$. Therefore
\begin{equation}
 \TV\left(b*\Law\left(\sqrt{3\delta/n}\sum_{j=1}^n U_j\right),
               b*\gamma_\delta\right)
 \sim\frac{\delta^2}{40n}\norm{(b*\gamma_\delta)^{(4)}}_1.
 \label{eq:equaldust}
\end{equation}
A fixed collection of macroscopic factors, or a fixed additional Gaussian,
may be absorbed into $b$; its derivative norms can only decrease under such
convolution. This gives a quantitative version of the perturbations used in
\cref{thm:continuity}.

\section{Exact finite-sample impossibility without a tail condition}
\label{sec:minimax}

Let $Y_1,\ldots,Y_N$ be independent with common law $P_\theta$,
$\theta=(v,a)\in\K$. Estimators may be any measurable functions of the data;
randomization does not improve the lower bounds below.

\begin{theorem}[No uniform statistical gain at any finite sample size]
\label{thm:minimax}
For every integer $N\ge1$,
\begin{align}
 \inf_{\widehat s}\sup_{\theta\in\K}
       \E_\theta|\widehat s-s(\theta)|&=\frac V2,
 \label{eq:absminimax}\\
 \inf_{\widehat s}\sup_{\theta\in\K}
       \E_\theta(\widehat s-s(\theta))^2&=\frac{V^2}4.
 \label{eq:sqminimax}
\end{align}
Both upper bounds are attained by the constant estimator $\widehat s=V/2$.
The lower bounds hold even on the fixed-total-variance slice $v=V$.
\end{theorem}
\begin{proof}
Compare the parameter $\theta_G$ with no uniform factors and Gaussian variance
$V$ to $\theta_n$ with zero Gaussian variance and $n$ equal half-lengths
$\sqrt{3V/n}$. Both have total variance $V$.
By \cref{prop:compact}, or by \eqref{eq:equaldust},
$\TV(P_{\theta_n},P_{\theta_G})\to0$. A product coupling, or telescoping
product densities, gives
\[
 \TV(P_{\theta_n}^{\otimes N},P_{\theta_G}^{\otimes N})
 \le N\TV(P_{\theta_n},P_{\theta_G})\longrightarrow0
\]
for every fixed $N$.

Clipping an estimator into $[0,V]$ cannot increase either loss. Let $f_n,f_G$
be the two sample densities, and let $z$ be the clipped estimate at a given
sample. The sum of the absolute-error risks is at least
\[
 \int\min(f_n,f_G)\bigl(|z|+|z-V|\bigr)
 =V\bigl(1-\TV(P_{\theta_n}^{\otimes N},P_{\theta_G}^{\otimes N})\bigr).
\]
The larger risk is at least half this quantity. Sending $n\to\infty$ gives
$V/2$. For squared error, use
$z^2+(z-V)^2\ge V^2/2$ in the same calculation to obtain $V^2/4$.
The constant estimator gives the matching upper bounds.
\end{proof}

\begin{corollary}[The corresponding testing obstruction]
\label{cor:test}
In testing the single law with $s=V,a=0$ against all laws with $s=0,v=V$,
the infimum, over tests, of type-I error plus worst-case type-II error equals
one at every finite sample size.
\end{corollary}
\begin{proof}
The sum of the errors for a simple pair is at least one minus their TV
distance. Use the preceding sequence of alternatives. A constant test attains
one.
\end{proof}

\begin{remark}[What this does and does not say]
Exact identification from an entire distribution is compatible with
\cref{thm:minimax}. The worst parameter is allowed to depend on the sample
size, and can contain as many tiny factors as needed. This is not a claim
that every fixed parameter is statistically indistinguishable from every
other one, or that pointwise consistent estimation is impossible. It also
does not contradict the repository's fixed-capacity rates: the required
number of factors here has no a priori bound.
\end{remark}

\begin{corollary}[Local minimax risk at an individual law]
\label{cor:localminimax}
Fix $\theta\in\K$ with Gaussian variance $s_0$. For $\eta>0$, restrict the
minimax problem to parameters $\theta'$ with
$\TV(P_{\theta'},P_\theta)<\eta$. For each fixed sample size $N$, as
$\eta\downarrow0$ the absolute-error minimax risk tends to $s_0/2$, and
the squared-error minimax risk tends to $s_0^2/4$.
\end{corollary}
\begin{proof}
Replace all Gaussian variance $s_0$ at $\theta$ by equal dust, keeping its
original uniforms. This produces parameters of Gaussian variance zero whose
observed laws converge in TV to $P_\theta$. Eventually they lie in any fixed
neighborhood, and the two-point proof of \cref{thm:minimax} yields lower
bounds $s_0/2$ and $s_0^2/4$. Upper semicontinuity of $s$, transferred to the
observed-law space by \cref{prop:compact}, implies that for every
$\varepsilon>0$ all sufficiently small neighborhoods have
$0\le s(\theta')\le s_0+\varepsilon$. The appropriate midpoint constant
estimator gives upper bounds $(s_0+\varepsilon)/2$ and
$(s_0+\varepsilon)^2/4$. Send $\varepsilon$ to zero.
\end{proof}

\section{Positive finite-cumulant certificates}
\label{sec:ladder}

A polynomial that equals one at zero but suppresses positive spectral atoms
provides an upper bound on the Gaussian variance. The simplest such family
already resolves the qualitative recovery question.

\begin{theorem}[A monotone cumulant ladder]
\label{thm:ladder}
For $m\ge0$, define
\begin{equation}
 A_m=v-\frac13\sum_{r=1}^m(-1)^{r+1}\binom mr\frac{p_{r+1}}{L^r},
 \qquad A_0=v.
 \label{eq:ladder}
\end{equation}
These quantities depend on cumulants only through order $2m+2$. They satisfy
\begin{equation}
 A_m-s=\frac13\sum_{j\ge1}x_j(1-x_j/L)^m\ge0,
 \qquad A_m\downarrow s.
 \label{eq:ladderbias}
\end{equation}
For every $J,m\ge1$,
\begin{equation}
 3(A_m-s)\le\frac{JL}{em}+\sum_{j>J}x_j.
 \label{eq:laddertail}
\end{equation}
Conversely,
\begin{equation}
 \sum_{j>2m}x_j\le6(A_m-s),\qquad m\ge1.
 \label{eq:ladderconverse}
\end{equation}
\end{theorem}
\begin{proof}
Integrate $(1-x/L)^m$ against \eqref{eq:spectral}, expand by the binomial
formula, and divide by three. The atom at zero contributes $s$.
For each positive $x_j$, $(1-x_j/L)^m\downarrow0$, and the summable
majorant $x_j$ permits dominated convergence.
For \eqref{eq:laddertail}, bound the first $J$ terms by
$x e^{-mx/L}\le L/(em)$ each, and keep the remaining terms unchanged.
For $j>2m$, $x_j\le L/j\le L/(2m+1)$, so Bernoulli's inequality gives
$(1-x_j/L)^m\ge1-m/(2m+1)>1/2$. Substitute in
\eqref{eq:ladderbias}.
\end{proof}

Every $A_m$ is continuous on $\K$. Indeed $v$ is continuous, and for $k\ge2$
the tail of $p_k$ is bounded by $L^k/(J+1)^{k-1}$, uniformly over $\K$.
The unknown $p_1$ never appears in \eqref{eq:ladder}; it would be circular
to use it as input to an estimator of $s$.

\section{A necessary and sufficient condition for uniform recovery}
\label{sec:criterion}

\begin{theorem}[Complete compact-class criterion]
\label{thm:criterion}
Let $K$ be a nonempty compact subset of $\K$, in the metric $d_\star$.
The following statements are equivalent:
\begin{enumerate}[label=\textup{(\roman*)}]
\item The Gaussian variance $s$ is continuous on $K$.
\item The uniform-factor energy tails vanish uniformly:
\begin{equation}
 \lim_{J\to\infty}\sup_{(v,a)\in K}\sum_{j>J}a_j^2=0.
 \label{eq:tailcriterion}
\end{equation}
\item The finite-cumulant ladder $A_m$ converges uniformly to $s$ on $K$.
\item There exist estimators from $N$ independent observations with
\begin{equation}
 \sup_{\theta\in K}\E_\theta|\widehat s_N-s(\theta)|\longrightarrow0.
 \label{eq:uniformconsistent}
\end{equation}
\end{enumerate}
\end{theorem}
\begin{proof}
For (i)$\Rightarrow$(ii), $p_1=3(v-s)$ is continuous on $K$ and its finite
partial sums are continuous and increase pointwise to $p_1$. Dini's theorem
gives uniform convergence. For (ii)$\Rightarrow$(i), the uniform limit of
these continuous partial sums is continuous, hence so is $s$.

The implication (ii)$\Rightarrow$(iii) follows by first choosing $J$ in
\eqref{eq:laddertail} and then $m$. Conversely \eqref{eq:ladderconverse}
proves (iii)$\Rightarrow$(ii). Alternatively, (i)$\Rightarrow$(iii) follows
from Dini's theorem applied to $A_m\downarrow s$.

For (iii)$\Rightarrow$(iv), fix $m$. Estimate the finitely many raw moments
up to order $2m+2$ by sample averages and clip them to fixed compact intervals
containing all the corresponding true moments on $\K$. Uniform moment
bounds from \cref{lem:subgaussian} make the expected moment errors
$O_m(N^{-1/2})$. Cumulants are polynomials in these moments, so their
errors have the same bound after clipping; therefore the estimator of
$A_m$ obtained from \eqref{eq:ladder} has uniformly vanishing expected error.
Choose increasing integers $N_m$ so large that this error is at most $1/m$
for all $N\ge N_m$. Let $m(N)$ increase sufficiently slowly along these
thresholds. Uniform convergence of $A_m-s$ and clipping the final answer
into $[0,V]$ prove \eqref{eq:uniformconsistent}.

Finally assume (iv), and clip the estimators into $[0,V]$. For each fixed $N$,
$\theta\mapsto\E_\theta\widehat s_N$ is continuous on $K$:
\cref{prop:compact} gives TV continuity of the one-observation laws,
tensorization gives it for their $N$-fold products, and expectations of a
bounded statistic are TV-continuous. By (iv), these continuous functions
converge uniformly to $s$. Hence $s$ is continuous.
\end{proof}

\begin{corollary}[Useful sufficient classes]
\label{cor:classes}
Uniform consistency holds on every compact subclass with any one of the
following bounds, with constants independent of the parameter: a fixed finite
capacity; $a_j\le Cj^{-\beta}$ for $\beta>1/2$; $a_j\le Cq^j$ for $0<q<1$;
or $\sum_j a_j\le A$.
\end{corollary}
\begin{proof}
The first three bounds give a summable deterministic majorant for $a_j^2$.
For the last, monotonicity gives $a_j\le A/j$, so
$\sum_{j>J}a_j^2\le a_{J+1}\sum_{j>J}a_j\le A^2/(J+1)$.
Apply \cref{thm:criterion}.
\end{proof}

\begin{remark}[Compactness and quantifiers matter]
An arbitrary nonclosed parameter class need not satisfy this equivalence:
continuity at each of its points does not imply a uniform modulus.
The theorem requires compactness in the specified topology, not merely a
bounded list of individual variances. The full $\K$ violates the tail
condition: for any fixed $J$, a sufficiently long equal-factor cloud can
place almost all of its energy beyond $J$.
\end{remark}

\section{Christoffel recovery from higher cumulants}
\label{sec:christoffel}

The binomial ladder is explicit, but it uses a prescribed polynomial.
Optimizing the polynomial yields the classical Christoffel function at zero,
applied here to the spectral measure rather than to the observed density.
This distinction is essential: empirical moments of $Y$ are not the moments
of $\rho$ until they have been converted by \cref{lem:cumulants}.

\begin{definition}[Gaussian-variance Christoffel certificate]
For an integer $d\ge0$, put
\begin{equation}
 \Cs_d(v,a)=\frac13
   \inf\left\{\int q(x)^2\dd\rho_{v,a}(x):
          q\in\R[x],\ \deg q\le d,\ q(0)=1\right\}.
 \label{eq:christoffel}
\end{equation}
\end{definition}

\begin{theorem}[Hierarchy, matrix formula, and tail certificate]
\label{thm:christoffel}
The quantities $\Cs_d$ depend on observed cumulants only through order $4d+2$.
They decrease to $s$ and satisfy
\begin{equation}
 0\le\Cs_d-s\le\frac13\sum_{j>d}x_j.
 \label{eq:christoffeltail}
\end{equation}
For $d\ge1$, define
\begin{equation}
 H_d=(p_{i+j+1})_{i,j=1}^d,\qquad y_d=(p_{i+1})_{i=1}^d.
 \label{eq:hankel}
\end{equation}
Then $y_d$ belongs to the range of the positive semidefinite matrix $H_d$, and
\begin{equation}
 \Cs_d=v-\frac13 y_d^{\mathsf T}H_d^\dagger y_d,
 \label{eq:schur}
\end{equation}
where $H_d^\dagger$ is the Moore--Penrose inverse; $\Cs_0=v$.
Moreover $\Cs_d=s$ if and only if the spectrum has at most $d$ distinct
positive squared scales.
\end{theorem}
\begin{proof}
Write $\mu_a=\sum_j x_j\delta_{x_j}$. Since $q(0)=1$,
\begin{equation}
 \Cs_d=s+\frac13\lambda_d,\qquad
 \lambda_d=\inf_{\deg q\le d,\ q(0)=1}\int q^2\dd\mu_a.
 \label{eq:lambda}
\end{equation}
For $q(x)=1+\sum_{i=1}^d c_ix^i$, the integral against $\rho$ is
\begin{equation}
 3v+2c^{\mathsf T}y_d+c^{\mathsf T}H_dc.
 \label{eq:quadratic}
\end{equation}
This uses moments $m_0,\ldots,m_{2d}$ and hence $v,p_2,\ldots,p_{2d+1}$,
not the unknown $p_1$.
If $u\in\ker H_d$, the polynomial $\sum_i u_ix^i$ has zero squared integral
against $\mu_a$, so it vanishes $\mu_a$-almost everywhere. Its integral
against $\mu_a$ is therefore $u^{\mathsf T}y_d=0$. Thus $y_d$ is orthogonal
to $\ker H_d$ and lies in its range. Completing the square on that range
shows that the minimum is attained and proves \eqref{eq:schur}.

If $x_1,\ldots,x_d$ are positive, use the test polynomial
\begin{equation}
 q_d(x)=\prod_{i=1}^d(1-x/x_i).
 \label{eq:annihilator}
\end{equation}
It vanishes at the first $d$ entries, and for every remaining $x_j$ each
factor lies in $[0,1]$. Its squared integral is at most
$\sum_{j>d}x_j$. If there are fewer positive entries, use their finite product
instead. This proves \eqref{eq:christoffeltail}; summability makes its
right-hand side tend to zero. Monotonicity follows by increasing the
admissible polynomial space.

If there are at most $d$ distinct positive squared scales, a polynomial
vanishing at each of them gives $\lambda_d=0$. Conversely, the attained
minimum can be zero only if an admissible polynomial vanishes at every
positive atom. It has at most $d$ distinct real roots. This proves the final
assertion, including $d=0$.
\end{proof}

\begin{remark}[A small-order formula]
Provided $p_3>0$,
\begin{equation}
 \Cs_1=v-\frac{p_2^2}{3p_3}.
 \label{eq:firstcertificate}
\end{equation}
It is exact for any number of equal positive half-lengths, because only one
distinct squared scale is present. At the zero spectrum it is interpreted by
the pseudoinverse formula and equals $v=s$. The numerator and denominator
are recovered from the fourth and sixth cumulants. This formula should not
be implemented as an unregularized ratio when $p_3$ is small.
\end{remark}

\begin{remark}[What the optimization relaxes]
The Christoffel problem is a positive-measure polynomial problem. It does not
impose the extra integer-multiplicity restriction
$\rho(\{x\})/x\in\N$ at positive atoms. Therefore its upper certificate is
not asserted to be the largest Gaussian variance among all uniform-factor
models matching a finite list of cumulants. That sharper constrained inverse
problem remains separate.
\end{remark}

\section{An exact geometric formula}
\label{sec:geometric}

Suppose the unknown squared scales form an exact geometric sequence,
\begin{equation}
 x_j=Cr^j,\qquad j=0,1,2,\ldots,
 \qquad C>0,\quad0<r<1.
 \label{eq:geometric}
\end{equation}
The reindexing from $j=1$ elsewhere is only for convenience. Here
$p_k=C^k/(1-r^k)$ and the uniform variance is $C/[3(1-r)]$.

\begin{theorem}[Exact finite-degree Christoffel bias]
\label{thm:geometric}
For the spectrum \eqref{eq:geometric}, at every integer $d\ge0$,
\begin{equation}
 \boxed{\displaystyle
  \Cs_d-s=\frac{C(1-r)r^d}{3(1-r^{d+1})^2}.}
 \label{eq:geometricbias}
\end{equation}
In particular, $\Cs_d-s\sim C(1-r)r^d/3$ as $d\to\infty$.
\end{theorem}
\begin{proof}
All finite moment matrices of $\mu_a$ are positive definite because its
support is infinite. Completing the square in the constant coefficient gives
\begin{equation}
 \lambda_d=
 \frac{\det\left(C^{i+j+1}/(1-r^{i+j+1})\right)_{i,j=0}^d}
      {\det\left(C^{i+j+1}/(1-r^{i+j+1})\right)_{i,j=1}^d},
 \label{eq:detratio}
\end{equation}
with the empty denominator determinant equal to one.

We use the elementary Cauchy identity
\begin{equation}
 \det\left(\frac1{1-u_iv_j}\right)_{i,j=0}^d
 =\frac{\prod_{0\le i<k\le d}(u_i-u_k)
          \prod_{0\le j<\ell\le d}(v_j-v_\ell)}
        {\prod_{i,j=0}^d(1-u_iv_j)}.
 \label{eq:cauchy}
\end{equation}
One proof applies the ordinary Cauchy determinant to
$1/(u_i^{-1}-v_j)$ and divides row $i$ by $u_i$; the powers of the $u_i$
cancel. The ordinary identity follows by multiplying the determinant by
its common denominator: the numerator is alternating in each variable list,
hence divisible by both Vandermonde products; comparison at a simple pole,
inductively on the matrix size, gives coefficient one. This argument proves
\eqref{eq:cauchy} for the nonzero distinct variables used here, with all
denominators nonzero.

In \eqref{eq:detratio}, row and column factors of $C$ leave a single factor
$C$ in the quotient. Apply \eqref{eq:cauchy} with $u_i=r^i$ and
$v_j=r^{j+1}$. The factors introduced by row and column zero in the numerator
are
\[
 \prod_{i=1}^d(1-r^i)\prod_{j=1}^d(r-r^{j+1})
 =r^d\prod_{i=1}^d(1-r^i)^2.
\]
The introduced denominator factors are
\[
 \prod_{j=0}^d(1-r^{j+1})\prod_{i=1}^d(1-r^{i+1})
 =(1-r)\prod_{i=2}^{d+1}(1-r^i)^2.
\]
Their quotient is $r^d(1-r)/(1-r^{d+1})^2$.
Multiply by $C$ and use \eqref{eq:lambda} to prove the theorem.
\end{proof}

\subsection{The dyadic Fabius--Rvachev specialization}
For half-lengths $a_j=2^{-j}$, $j\ge1$, set $C=r=1/4$. The exact result is
\begin{equation}
 \Cs_d-s=\frac{4^{-d}}{16(1-4^{-d-1})^2}.
 \label{eq:dyadicbias}
\end{equation}
The first few values are particularly simple:
\begin{center}
\begin{tabular}{@{}rrrr@{}}
\toprule
$d$ & Highest cumulant order & Exact bias & Decimal bias \\\midrule
0 & 2  & $1/9$          & $0.1111111111$\\
1 & 6  & $4/225$        & $0.0177777778$\\
2 & 10 & $16/3969$      & $0.0040312421$\\
3 & 14 & $64/65025$     & $0.0009842368$\\
4 & 18 & $256/1046529$  & $0.0002446182$\\
\bottomrule
\end{tabular}
\end{center}
The accompanying program regenerates these rational values and the decimal
roundings. Unlike a finite truncation of the uniform sum, these numbers are
exact for the infinite geometric spectrum.

\subsection{Interpretation and its limitations}
In the zero-based indexing of \eqref{eq:geometric}, deleting the first $d$
entries leaves the tail $\sum_{j=d}^\infty Cr^j=Cr^d/(1-r)$.
Consequently \eqref{eq:geometricbias} improves the generic upper bound
in \eqref{eq:christoffeltail} by the exact factor
\begin{equation}
 \frac{\Cs_d-s}{Cr^d/[3(1-r)]}
 =\frac{(1-r)^2}{(1-r^{d+1})^2}\longrightarrow(1-r)^2.
 \label{eq:tailratio}
\end{equation}

This exact formula is not a statistical minimax rate. It describes noiseless
finite-moment bias. If exact geometric shape with unknown $C,r$ were imposed
as a parametric assumption, those few parameters could instead be recovered
from a few cumulants. The general-purpose certificate deliberately does not
use that parametric shortcut.

\section{Regularization by bounded polynomial coefficients}
\label{sec:regularization}

Direct inversion of $H_d$ is unnecessary for a robust upper certificate.
Normalize spectral locations by $z=x/L$ and define
\begin{equation}
 \bar m_0=3v,\qquad
 \bar m_r=p_{r+1}/L^r\quad(r\ge1).
 \label{eq:normalizedmoments}
\end{equation}
These are the moments of the pushforward of $\rho$ to $[0,1]$; the total
mass is still $3v$, not one.
For $M\ge1$, let
\begin{equation}
 \mathcal Q_{d,M}=
 \left\{q(z)=\sum_{i=0}^d c_iz^i:
 c_0=1,\ \sum_{i=0}^d|c_i|\le M\right\}.
 \label{eq:coeffball}
\end{equation}
This is a compact coefficient set.
Given estimates $\widehat m_0,\ldots,\widehat m_{2d}$, define
\begin{equation}
 \widehat C_{d,M}=\frac13\inf_{q\in\mathcal Q_{d,M}}
       \sum_{i,j=0}^d c_ic_j\widehat m_{i+j}.
 \label{eq:regularized}
\end{equation}
The value is a continuous function of the estimated moments, even if their
Hankel matrix is not positive semidefinite. The estimated quadratic may then
be nonconvex; existence of its global minimum still follows from compactness.
No efficient general nonconvex solver is implicit in this definition.

\begin{theorem}[A deterministic noise certificate]
\label{thm:regularized}
Suppose
\begin{equation}
 \max_{0\le r\le2d}|\widehat m_r-\bar m_r|\le\varepsilon.
 \label{eq:momenterror}
\end{equation}
Let $C_{d,M}$ be \eqref{eq:regularized} with the true moments.
Then
\begin{equation}
 |\widehat C_{d,M}-C_{d,M}|\le\frac{\varepsilon M^2}{3}.
 \label{eq:noisebound}
\end{equation}
If $x_i\ge\ell_i>0$ for $1\le i\le d$ and
\begin{equation}
 M\ge M_d:=\prod_{i=1}^d(1+L/\ell_i),
 \label{eq:Md}
\end{equation}
then
\begin{equation}
 |\clip(\widehat C_{d,M})-s|
 \le\frac13\sum_{j>d}x_j+\frac{\varepsilon M^2}{3}.
 \label{eq:regularizedrisk}
\end{equation}
Moreover
\begin{equation}
 U_{d,M}:=\clip\left(\widehat C_{d,M}+\frac{\varepsilon M^2}{3}\right)
 \label{eq:uppercertificate}
\end{equation}
is an upper bound on $s$, and
$0\le U_{d,M}-s\le\frac13\sum_{j>d}x_j+2\varepsilon M^2/3$.
\end{theorem}
\begin{proof}
For any $q\in\mathcal Q_{d,M}$, the difference between the estimated and true
quadratic is at most
$\varepsilon\sum_{i,j}|c_i||c_j|\le\varepsilon M^2$.
The infima differ by at most the same amount, proving \eqref{eq:noisebound}.
The polynomial $\prod_{i=1}^d(1-Lz/x_i)$ has coefficient $\ell^1$ norm
exactly $\prod_{i=1}^d(1+L/x_i)\le M_d$, because its coefficients alternate
in sign. It is therefore admissible. The argument of
\cref{thm:christoffel} gives
$0\le C_{d,M}-s\le\frac13\sum_{j>d}x_j$.
Combine this with \eqref{eq:noisebound}. Clipping into an interval containing
$s$ does not increase absolute error, and preserves the certified upper-bound
property.
\end{proof}

\begin{remark}[A fully specified, though possibly expensive, computation]
One may approximate the infimum in \eqref{eq:regularized} over a finite
rational grid in the coefficient ball. If the coefficient $\ell^1$ mesh is
$\eta$, and $K_0=\max_r|\widehat m_r|$, the objective changes by at most
$2K_0M\eta$ between two grid-neighbor coefficient vectors, since
$|c^{\mathsf T}\widehat Hc-d^{\mathsf T}\widehat Hd|
\le K_0\norm{c-d}_1(\norm c_1+\norm d_1)$.
A grid fine enough to meet a specified objective tolerance therefore gives
a finite globally certified procedure. Such a search is not promised to be
practical in high degree; the theorem is an error guarantee, not a complexity
claim.
\end{remark}

\section{Statistical upper bounds on envelope classes}
\label{sec:rates}

We now turn the deterministic certificate into actual estimation guarantees.
These bounds do not claim optimal rates. Their purpose is to prove a
constructive contrast with \cref{thm:minimax} and quantify the price of
infinite rank under explicit assumptions.

\begin{lemma}[Uniform estimation of the normalized spectral moments]
\label{lem:empirical}
For each $d\ge1$, there are estimates of
$\bar m_0,\ldots,\bar m_{2d}$ from $N$ observations such that
\begin{equation}
 \sup_{\theta\in\K}\E_\theta
   \max_{0\le r\le2d}|\widehat m_r-\bar m_r|
 \le N^{-1/2}\exp\{C_0d\log(d+1)\},
 \label{eq:empirical}
\end{equation}
where $C_0$ depends only on $H,V$ and the known background, not on $d,N$.
\end{lemma}
\begin{proof}
Set $R=4d+2$. Increase a constant $D_1$ so that, with
$M_R=D_1\sqrt R\ge1$, every raw moment $\mu_j=\E Y^j$, $j\le R$, has
$|\mu_j|\le M_R^j$ and $\E|Y|^{2j}\le M_R^{2j}$.
This follows from \cref{lem:subgaussian}, using $2j\le2R$ and absorbing
$\sqrt2$ in $D_1$. Estimate $\mu_j$ by the sample mean of $Y^j$ and clip
it into $[-M_R^j,M_R^j]$. The resulting $\widehat\mu_j$ satisfies
\begin{equation}
 \E|\widehat\mu_j-\mu_j|\le M_R^j/\sqrt N.
 \label{eq:rawmomenterror}
\end{equation}

For $r\le R$, use the standard partition expression for a cumulant,
\begin{equation}
 \kappa_r=\sum_{\pi\in\Pi_r}(-1)^{|\pi|-1}(|\pi|-1)!
              \prod_{A\in\pi}\mu_{|A|},
 \label{eq:partitioncumulant}
\end{equation}
which follows by comparing coefficients in
$\log(\sum_{j\ge0}\mu_jt^j/j!)$.
A telescoping product and \eqref{eq:rawmomenterror} bound the expected error
of the product for a partition $\pi$ by $|\pi|M_R^r/\sqrt N$.
Furthermore,
\[
 \sum_{\pi\in\Pi_r}(|\pi|-1)!|\pi|
 =\sum_{q=1}^r S(r,q)q!\le\sum_{q=1}^r q^r\le r^{r+1},
\]
because $S(r,q)q!$ counts surjections to $q$ labeled points. Thus the
expected cumulant error is at most $r^{r+1}M_R^r/\sqrt N$.

Subtract the exact known cumulants of $B$. Use \eqref{eq:variance} for
$\bar m_0=3v$ and \eqref{eq:cumulants} for $\bar m_r$, $r\ge1$.
By \eqref{eq:gamma}, $|\gamma_{2k}|^{-1}\le\pi^{2k}/2$.
The factors $L^{-r}$ and these reciprocal cumulants grow at most
exponentially in $R$, since $V>0$ is fixed. Sum the expected errors to bound
the expected maximum. All factors are absorbed in
$\exp\{C_0d\log(d+1)\}$, proving the lemma.
\end{proof}

\begin{theorem}[Constructive risk bound]
\label{thm:statbound}
Suppose a class satisfies $x_i\ge\ell_i>0$ for $i\le d$ and
$\sum_{j>d}x_j\le T_d$ uniformly. Using the estimates in
\cref{lem:empirical}, the clipped estimator \eqref{eq:regularized} with
$M=M_d$ satisfies
\begin{equation}
 \sup_\theta\E_\theta|\clip(\widehat C_{d,M_d})-s|
 \le\frac{T_d}{3}
 +\frac{M_d^2}{3\sqrt N}\exp\{C_0d\log(d+1)\}.
 \label{eq:statrisk}
\end{equation}
A finite coefficient-grid minimization with objective tolerance $N^{-1}$
adds at most $1/(3N)$ to this bound.
\end{theorem}
\begin{proof}
The deterministic argument of \cref{thm:regularized} holds for every realized
moment error, not only on a fixed high-probability event. Take expectations
and apply \cref{lem:empirical}. The grid tolerance statement follows directly
from division of the objective value by three.
\end{proof}

\begin{corollary}[Polynomial envelopes]
\label{cor:polyrate}
Suppose the sorted half-lengths satisfy
\begin{equation}
 c_-j^{-\beta}\le a_j\le c_+j^{-\beta}\quad(j\ge1),
 \qquad c_->0,\quad \beta>1/2,
 \label{eq:polyenvelope}
\end{equation}
with fixed constants, and the class is nonempty under the variance budget.
There exist estimators for which, as $N\to\infty$,
\begin{equation}
 \sup_\theta\E_\theta|\widehat s_N-s|
 =O\left(\left(\frac{\log N}{\log\log N}\right)^{1-2\beta}\right).
 \label{eq:polyrate}
\end{equation}
\end{corollary}
\begin{proof}
Take $\ell_i=c_-^2i^{-2\beta}$. Then
\[
 \log M_d\le d\log(1+L/c_-^2)+2\beta\log(d!)=O(d\log(d+1)),
 \qquad
 T_d\le\frac{c_+^2}{2\beta-1}d^{1-2\beta}.
\]
Choose $d=\lfloor c\log N/\log\log N\rfloor$ with $c>0$ sufficiently
small relative to the constants in \eqref{eq:statrisk}. Its noise term is
$O(N^{-1/4})$ for all sufficiently large $N$; its bias term is the displayed
bound. Handle finitely many small $N$ by the constant estimator.
\end{proof}

\begin{corollary}[Geometric envelopes]
\label{cor:geomrate}
Suppose
\begin{equation}
 c_-q^j\le a_j\le c_+q^j\quad(j\ge1),
 \qquad c_->0,\quad0<q<1,
 \label{eq:geomenvelope}
\end{equation}
with fixed constants and nonempty model class. There are constants $c,C>0$
and estimators such that
\begin{equation}
 \sup_\theta\E_\theta|\widehat s_N-s|
 \le C\exp\{-c\sqrt{\log N}\}
 \label{eq:geomrate}
\end{equation}
for all sufficiently large $N$.
\end{corollary}
\begin{proof}
Take $\ell_i=c_-^2q^{2i}$. Then
\[
 \log M_d\le d\log(1+L/c_-^2)+\log(1/q)d(d+1)=O(d^2),
 \qquad
 T_d\le\frac{c_+^2q^{2d+2}}{1-q^2}.
\]
Set $d=\lfloor c_1\sqrt{\log N}\rfloor$ with $c_1$ sufficiently small.
The noise term in \eqref{eq:statrisk} is $O(N^{-1/4})$ for large $N$, and
the tail term is bounded by the right-hand side of \eqref{eq:geomrate}.
\end{proof}

\begin{remark}[Precisely what assumptions these rates use]
The positive lower envelopes control the coefficients of the annihilating
polynomial. They are not needed for the qualitative uniform-consistency
criterion. Neither rate is asserted minimax, and neither is a rate for
estimating the entire ordered spectrum. These classes allow independent
deviations from geometric or polynomial shape at every coordinate; they are
not finite-dimensional parametric families. The constants depend on the fixed
envelope and may deteriorate badly as its lower bound decreases.
\end{remark}

\section{Pointwise estimation despite constant global minimax risk}
\label{sec:pointwise}

The global obstruction concerns a supremum over parameters; it should not be
mistaken for an impossibility of learning any single fixed distribution.
The positive ladder produces one sequence of estimators that is consistent
at every fixed member of the full unrestricted class.

\begin{theorem}[One pointwise consistent estimator on the whole class]
\label{thm:pointwise}
Let $m=m(N)=\lfloor\sqrt{\log(N+e)}\rfloor$. Use
\cref{lem:empirical} with $d=m$ to estimate the normalized moments, and put
\begin{equation}
 \widehat s_N=
 \clip\left(\frac13\sum_{r=0}^{m}(-1)^r\binom mr\widehat m_r\right).
 \label{eq:pointwiseestimator}
\end{equation}
For every fixed $\theta\in\K$,
$\E_\theta|\widehat s_N-s(\theta)|\to0$.
More precisely, uniformly over $\theta$,
\begin{equation}
 \E_\theta|\widehat s_N-s(\theta)|
 \le A_m(\theta)-s(\theta)
  +\frac{2^m}{3\sqrt N}\exp\{C_0m\log(m+1)\}.
 \label{eq:pointwiserisk}
\end{equation}
Only the first term need fail to converge uniformly.
\end{theorem}
\begin{proof}
The true moment polynomial in \eqref{eq:pointwiseestimator} is $A_m$.
The sum of the absolute binomial coefficients is $2^m$, so
\cref{lem:empirical} bounds the expected error of this polynomial by the
second term of \eqref{eq:pointwiserisk}. Clipping cannot increase the error
against $s$. The second term tends to zero since
$m\log(m+1)=o(\log N)$. For each fixed parameter, \cref{thm:ladder} makes
the first term tend to zero as well.
\end{proof}

The observations therefore support the following three simultaneous facts:
exact-law identifiability; pointwise consistent estimation from increasing
samples; and a nonzero, sample-size-independent global minimax risk. The
resolution is the absence of a uniform energy-tail modulus. The dust
alternative in a worst-case lower bound becomes increasingly complicated as
the sample size grows.

\section{Verification, reproducibility, and formalization route}
\label{sec:validation}

\subsection{What the accompanying program checks}
The file \path{verify_results.py} uses exact rational arithmetic for algebraic
checks and high-precision arithmetic for a separately labeled Fourier
diagnostic. It checks the Bernoulli cumulants, all constants in the Taylor
remainder, the geometric Schur-complement formula over several rational
values of $C,r$, finite-spectrum pseudoinverse formulas and exact termination,
positive-ladder identities, and representative coefficient-noise bounds.
It also regenerates the dyadic table. A high-precision calculation of equal-dust
characteristic functions tests the sign and coefficient of the first correction.
That diagnostic is \emph{not} a numerical computation of total variation.

The delivered run passed 4,619 exact assertions and three high-precision
Fourier diagnostic checks. All data written by the script are consequences or diagnostics of the stated
formulas. They are not used to bridge an unproved infinite range. In
particular, the statistical assertions are justified by the proofs, not by
Monte Carlo simulations. The program is an independently executable
companion, not a formal proof certificate accepted by Lean.

\subsection{A realistic formalization decomposition}
A future Lean or Rocq development can separate five layers.
First formalize the finite algebraic layer: moment-cumulant conversion,
Hankel Gram matrices, a quadratic minimum on a matrix range, and the Cauchy
determinant quotient. This would certify the geometric formula and finite
spectral computations without any probability theory.

Second formalize square-summable independent sums and their exponential
bounds, then the spectral measure and compactification. Third establish
convolution continuity in $L^1$ and the operator Taylor estimates; finite
products should precede the infinite limit. Fourth formalize the two-point
risk lower bound and the tensorization inequality. Finally add Dini's theorem,
empirical moments, and the diagonal estimator for the compact-class criterion.

No declaration in the current repository is asserted to implement these new
endpoints. The mathematical proof dependencies are explicit enough that one
can formalize the layers independently rather than importing the conclusions
as axioms.

\section{Further research questions and topics}
\label{sec:questions}

\begin{question}[Sharp minimax rates under spectral envelopes]
Determine matching lower and upper rates for Gaussian-variance estimation
under the two-sided envelopes \eqref{eq:polyenvelope} and
\eqref{eq:geomenvelope}. Are the logarithmic and stretched-logarithmic upper
bounds proved here optimal, or does the integer-multiplicity constraint
permit a substantially faster method?
\end{question}
The compactness criterion settles possibility, not speed. A lower bound must
preserve the fixed envelope while producing moment or likelihood cancellation;
the unrestricted equal-dust construction does not do this.

\begin{question}[Removing lower envelopes quantitatively]
Obtain sharp rates assuming only $a_j\le Cj^{-\beta}$ or $a_j\le Cq^j$.
Can positive polynomials adapted to a scale threshold replace the explicit
root polynomial without incurring its coefficient growth?
\end{question}
Uniform consistency already follows from \cref{thm:criterion}. What remains
is an effective quantitative modulus that does not presume a positive lower
bound for every leading squared scale.

\begin{question}[Optimal recovery from finitely many exact cumulants]
Given $v,p_2,\ldots,p_m$, characterize the largest and smallest possible
Gaussian variances among actual uniform-factor models. Determine when a
Christoffel upper bound is attainable subject to the discrete condition
$\rho(\{x\})/x\in\N$ for every positive atom.
\end{question}
Classical truncated moment feasibility alone drops this condition. This is a
natural place where the special convolution model may improve on a generic
positive-measure relaxation.

\begin{question}[A stable computational implementation]
Design a numerically efficient estimator with certified error bounds from
noisy cumulants. Compare coefficient balls, orthogonal-polynomial bases,
positive-semidefinite moment projection, and interval arithmetic. A useful
implementation should report a noise-dependent degree rather than simply
invert a large Hankel matrix.
\end{question}
The exact formulas in this article establish a reference target; the finite
rational-grid construction establishes existence, not practical scalability.

\begin{question}[Perturbations of geometric spectra]
For $x_j=Cr^j(1+\varepsilon_j)$ with controlled $\varepsilon_j$, determine
whether the exact geometric Christoffel asymptotic persists with a stable
leading constant. Identify the weakest summability or regular-variation
assumptions yielding a relative-error theorem.
\end{question}
A first route is determinant perturbation, but its conditioning may be much
worse than that of the variational problem. Direct comparison of positive
spectral measures may be more robust.

\begin{question}[Higher-order dust expansions]
Extend \cref{thm:dust} to arbitrary fixed order with constants independent of
rank. Express the remainder in a small collection of power sums and derive
uniform relative errors when several leading cumulants have been matched.
\end{question}
The double-telescoping proof suggests a route via higher operator differences.
The issue is controlling the remainder without an exponentially large subset
expansion. Existing geometric Edgeworth reports address a different scaling
regime and should be compared carefully before any novelty claim.

\begin{question}[The unsmoothed statistical experiment]
Which topology and local risk statements survive when the known smooth
background is removed, particularly at the point mass and at finite uniform
sums? Derive sharp distinctions between weak, total-variation, and Hellinger
behavior at support boundaries.
\end{question}
The smoothing assumption is used explicitly in the TV-homeomorphism proof.
It must not be silently discarded, even though some dust approximations
have stronger convergence without smoothing.

\begin{question}[Unknown or partly known background]
What aspects of Gaussian variance remain identifiable when $B$ is known only
through finitely many moments, or belongs to a specified compact family?
Formulate sufficient restrictions that distinguish background Gaussian
content from the latent Gaussian component.
\end{question}
The present subtraction of background cumulants is exact. With an arbitrary
unknown background, convolution factorization creates genuine nonidentifiability,
not merely statistical instability.

\begin{question}[Multivariate segment convolutions]
Replace $a_jU_j$ by $u_jU_j$ with $u_j\in\R^p$, and Gaussian variance by a
positive semidefinite covariance matrix. Characterize the corresponding
matrix-valued dust, directional cumulant measures, and compact-class
learnability conditions.
\end{question}
In one dimension the ordered nonnegative spectrum removes sign and permutation
ambiguity. In several dimensions, opposite directions and angular moment
identification introduce additional structure that must be resolved first.

\begin{question}[A proof-assistant endpoint with statistical semantics]
Formalize \cref{thm:criterion} together with \cref{thm:minimax} for the actual
probability-law model, not only for finite algebraic moment data. Can the
spectral-measure representation reduce the amount of special-function
analysis needed in that formalization?
\end{question}
The finite Schur-complement and determinant lemmas give a manageable initial
milestone, while the compact statistical experiment provides a substantial
longer-term target directly aligned with ProveIt's verification goals.

\section{Conclusion}
The important distinction is not simply finite versus infinite sums. It is
whether a class permits a fixed amount of variance to migrate into arbitrarily
many unresolved factors. In the full class, that migration produces exact
constant minimax risks at every finite sample size. On compact subclasses,
uniformly vanishing energy tails are exactly the condition that restores
uniform Gaussian-variance estimation.

The variance-weighted spectral measure makes this distinction transparent:
the Gaussian component is its mass at zero. Positive polynomial filters and
Christoffel minimization recover this atom from higher cumulants. The exact
geometric formula supplies a particularly explicit link back to the
Fabius--Rvachev setting, and bounded-coefficient regularization translates the
algebra into honest noise and sampling guarantees. The remaining open work is
chiefly quantitative: sharp envelope rates, efficient certified algorithms,
and extensions beyond the known smooth background.

\begin{thebibliography}{99}
\bibitem{repo}
Vladimir Reshetnikov, \emph{ProveIt}, GitHub repository, inspected at commit
\texttt{458ccfb80b9940220819091da107e82bb970dcca}, 29 September 2026.
\url{https://github.com/VladimirReshetnikov/ProveIt}.

\bibitem{finite}
\emph{Sharp Stability Strata for Finite Fabius--Rvachev Deconvolution:
Colliding scales, vanishing factors, and optimal recovery exponents},
AI-assisted research report prepared for Vladimir Reshetnikov,
28 September 2026; archived in \cite{repo}.
\url{https://github.com/VladimirReshetnikov/ProveIt/tree/458ccfb80b9940220819091da107e82bb970dcca/Analysis/FabiusFunction/docs/semi-formalized-research-frontiers/drafts/inverse-and-sampling/Sharp_Stability_Strata_Fabius_Rvachev_Deconvolution}.

\bibitem{gaussian}
\emph{Gaussian Confounding and Sharp Recovery of Uniform Convolution Factors:
Shifted power sums, minimax estimation, and detection above a Fabius--Rvachev
background}, AI-assisted research report prepared for Vladimir Reshetnikov,
29 September 2026; archived in \cite{repo}.
\url{https://github.com/VladimirReshetnikov/ProveIt/tree/458ccfb80b9940220819091da107e82bb970dcca/Analysis/FabiusFunction/docs/semi-formalized-research-frontiers/drafts/inverse-and-sampling/Gaussian_Confounding_Sharp_Recovery_Uniform_Factors}.

\bibitem{recovering}
\emph{Recovering Uniform Factors: Exact Moment Fibres and Logarithmic
Instability Near the Fabius--Rvachev Law}, AI-assisted research report,
archived in \cite{repo}, September 2026.
\url{https://github.com/VladimirReshetnikov/ProveIt/tree/458ccfb80b9940220819091da107e82bb970dcca/Analysis/FabiusFunction/docs/semi-formalized-research-frontiers/drafts/inverse-and-sampling/Recovering_Uniform_Factors_Fabius_Rvachev}.

\bibitem{native}
\emph{Fabius--Rvachev New Frontiers}, representation-frontier package,
\emph{Corpus audit and nonduplication boundary}, inspected in \cite{repo}.
\url{https://github.com/VladimirReshetnikov/ProveIt/blob/458ccfb80b9940220819091da107e82bb970dcca/Analysis/FabiusFunction/docs/semi-formalized-research-frontiers/drafts/representations/Fabius_Rvachev_New_Frontiers-2/CORPUS_AUDIT.md}.

\bibitem{billey}
Sara C. Billey and Joshua P. Swanson,
\emph{The metric space of limit laws for $q$-hook formulas},
arXiv:2010.12701, 2020; version 2, 24 August 2023.
\url{https://arxiv.org/abs/2010.12701v2}.

\bibitem{christoffel}
Edouard Pauwels, Mihai Putinar, and Jean-Bernard Lasserre,
\emph{Data analysis from empirical moments and the Christoffel function},
arXiv:1810.08480, 2018; version 2, 10 January 2020.
\url{https://arxiv.org/abs/1810.08480v2}.

\bibitem{hausdorff}
Daniel Gerth, Bernd Hofmann, Christopher Hofmann, and Stefan Kindermann,
\emph{The Hausdorff Moment Problem in the light of ill-posedness of type I},
arXiv:2104.06029, 2021.
\url{https://arxiv.org/abs/2104.06029}.

\bibitem{dlmf}
NIST Digital Library of Mathematical Functions,
\emph{Section 4.19: Maclaurin Series and Laurent Series},
especially Eq.~4.19.7; accessed 29 September 2026.
\url{https://dlmf.nist.gov/4.19}.

\bibitem{dlmfq}
NIST Digital Library of Mathematical Functions,
\emph{Section 18.27: $q$-Hahn Class}, especially the little $q$-Jacobi
orthogonality in Section 18.27(iv); accessed 29 September 2026.
\url{https://dlmf.nist.gov/18.27}.
\end{thebibliography}
\end{document}
